\documentclass[11pt,a4paper]{article}

% =============================================================================
% Paquetes
% =============================================================================
\usepackage[utf8]{inputenc}
\usepackage[T1]{fontenc}
\usepackage[spanish,es-minimal]{babel}
\usepackage{amsmath,amssymb,amsthm}
\usepackage{geometry}
\geometry{margin=2.6cm}
\usepackage{booktabs}
\usepackage{longtable}
\usepackage{array}
\usepackage{enumitem}
\usepackage{tikz}
\usepackage{float}
\usepackage{hyperref}
\hypersetup{
  colorlinks=true,
  linkcolor=blue!60!black,
  urlcolor=blue!60!black,
  citecolor=blue!60!black,
  pdftitle={El proceso de cálculo hidráulico 1D del motor LocalSWMM}
}
\usepackage{microtype}
\usepackage{xcolor}
\usepackage{siunitx}

% =============================================================================
% Entornos y macros
% =============================================================================
\newtheorem{definition}{Definición}[section]
\newtheorem{remark}{Observación}[section]

\newcommand{\st}{\mathrm{s}}
\newcommand{\ft}{\mathrm{ft}}
\newcommand{\cfs}{\mathrm{ft^3/s}}
\newcommand{\sfft}{\mathrm{ft^2}}
\newcommand{\R}{\mathbb{R}}
\renewcommand{\arraystretch}{1.25}

% Abreviatura para rutas del motor
\newcommand{\src}[1]{\texttt{src/engine/#1}}

\title{\textbf{El Proceso de Cálculo Hidráulico 1D}\\[0.6em]
\large del proyecto LocalSWMM (motor HydroCouple OpenSWMM)}
\author{Documento de Referencia Técnica}
\date{Agosto 2026}

\begin{document}

\maketitle

\begin{abstract}
Este documento entrega una descripción completa, a nivel de ecuaciones, del
proceso de cálculo hidráulico unidimensional (1D) que utiliza la aplicación web
LocalSWMM. Todos los cálculos hidráulicos los ejecuta el motor HydroCouple
\textbf{OpenSWMM}, compilado a WebAssembly y que se ejecuta dentro del
navegador. El motor resuelve las ecuaciones de de~Saint-Venant sobre una red de
nodos y enlaces. Se dispone de cuatro formulaciones de tránsito (routing):
\emph{escurrimiento permanente} (flujo permanente), \emph{onda cinemática},
\emph{onda dinámica} (la opción por defecto) y un esquema explícito de
\emph{volúmenes finitos}. El documento explica el modelo conceptual, las
ecuaciones de gobierno, los esquemas de diferencias finitas, el procedimiento
de solución iterativa, los algoritmos de continuidad en los nodos y de
escurrimiento en carga (surcharge), el cálculo de las estructuras hidráulicas
no-conducto, las funciones de geometría de sección transversal, las
condiciones de borde, el paso de tiempo adaptativo y el balance de masa.
Toda ecuación está transcrita del código fuente del motor (ver la lista de
archivos en las Referencias) y referenciada a los Manuales de Referencia del
motor cuando corresponde.
\end{abstract}

\tableofcontents
\newpage

% =============================================================================
% 1. INTRODUCCIÓN
% =============================================================================
\section{Introducción}

\subsection{Alcance}

Este documento cubre los cálculos hidráulicos (tránsito de caudales)
\emph{unidimensionales} del proyecto LocalSWMM. LocalSWMM es una aplicación web
basada en navegador para el modelamiento de redes de aguas lluvias y aguas
servidas; incorpora el motor HydroCouple OpenSWMM (C$++$) compilado a
WebAssembly. El componente hidráulico 1D resuelve el escurrimiento
unidimensional, no permanente y gradualmente variado a través de una red de
\emph{nodos} (cámaras de unión, estanques de almacenamiento, divisores de
flujo, emisarios) y \emph{enlaces} (conductos, bombas, orificios, vertederos,
descargas). El módulo de escurrimiento superficial 2D y su acoplamiento 1D--2D
quedan fuera del alcance, salvo donde interactúan con el solver 1D (por
ejemplo, a través de los aportes laterales y del acoplamiento por área de
encharcamiento).

El motor es un \emph{modelo de simulación distribuida de tiempo discreto}:
avanza un vector de estado a lo largo de una secuencia de pasos de tiempo,
\begin{equation}
  \boldsymbol{X}_{t} = f\!\left(\boldsymbol{X}_{t-1}, \boldsymbol{I}_{t}, \boldsymbol{P}\right),
  \qquad
  \boldsymbol{Y}_{t} = g\!\left(\boldsymbol{X}_{t}, \boldsymbol{P}\right),
\end{equation}
donde $\boldsymbol{X}$ es el vector de estado, $\boldsymbol{I}$ las entradas
externas, $\boldsymbol{P}$ los parámetros fijos e $\boldsymbol{Y}$ las salidas.

\subsection{Glosario de terminología (español de Chile)}

A lo largo de este documento se usa la terminología técnica empleada en la
práctica hidrosanitaria chilena. La Tabla~\ref{tab:glosario} entrega la
equivalencia entre los términos del motor (en inglés) y su traducción local.

\begin{longtable}{>{\raggedright}p{6.2cm}>{\raggedright\arraybackslash}p{9.2cm}}
\toprule
\textbf{Término del motor} & \textbf{Español de Chile} \\
\midrule
\endhead
stormwater & aguas lluvias \\
wastewater / sewage & aguas servidas \\
sewer system & alcantarillado \\
sewer main (collector) & colector \\
culvert & alcantarilla de paso \\
manhole / junction chamber & cámara de inspección \\
junction node & cámara de unión \\
outfall & emisario / punto de descarga \\
storage unit & estanque de almacenamiento (piscina de laminación / estanque de retención) \\
pump & bomba \\
pump station & cámara de bombeo / estación de bombeo \\
wet well & cámara húmeda \\
weir & vertedero \\
orifice & orificio \\
outlet (link) & descarga regulada / salida con curva de gasto \\
flap gate & compuerta de retención \\
link / conduit & vínculo / enlace; conducto / tubería \\
invert elevation & cota de batea \\
crown & clave \\
rim & brocal \\
flow depth & tirante \\
water surface width & espejo de agua / ancho superficial \\
hydraulic radius & radio hidráulico \\
hydraulic head & carga hidráulica / cota piezométrica \\
friction slope & pendiente de fricción / pendiente motriz \\
minor losses & pérdidas locales / pérdidas singulares \\
backwater & remanso \\
surcharge (state) & escurrimiento en carga \\
flowing full & a sección llena \\
overflow & rebalse \\
flooding & anegamiento / inundación \\
ponding & encharcamiento \\
runoff & escorrentía superficial \\
groundwater table & napa freática \\
infiltration & infiltración \\
exfiltration & exfiltración \\
street inlet / catch basin & sumidero \\
rain gage & pluviómetro \\
catchment / subcatchment & cuenca / subcuenca \\
time step & paso de tiempo \\
flow rate & caudal \\
hydrograph & hidrograma \\
time series & serie de tiempo \\
control rules & reglas de control \\
steady flow & escurrimiento permanente \\
kinematic wave & onda cinemática \\
dynamic wave & onda dinámica \\
\bottomrule
\caption{Glosario de terminología técnica en español de Chile.}
\label{tab:glosario}
\end{longtable}

\subsection{El modelo conceptual de nodos y enlaces}

La porción de conducción de un sistema de drenaje es una red de nodos y
enlaces. Los \textbf{nodos} son puntos que representan:
\begin{itemize}[nosep]
  \item \textbf{Cámaras de unión} --- puntos donde se unen los enlaces, con
        cota de batea, altura hasta el brocal, profundidad opcional de
        escurrimiento en carga (presión) y área opcional de encharcamiento.
  \item \textbf{Emisarios (puntos de descarga)} --- nodos de borde terminales
        con una condición de nivel de agua prescrita (libre, normal, fija, de
        marea o serie de tiempo) y compuerta de retención opcional.
  \item \textbf{Estanques de almacenamiento} --- nodos con una relación
        superficie--versus--tirante y volumen de almacenamiento real; nunca se
        presurizan.
  \item \textbf{Divisores de flujo} --- nodos que dividen el ingreso entre un
        enlace de continuación y un enlace de desvío (de corte, de rebalse,
        tabular o de vertedero; activos bajo onda cinemática, tratados como
        cámaras de unión bajo onda dinámica).
\end{itemize}

Los \textbf{enlaces} conectan nodos. Los tipos de enlace son:
\begin{itemize}[nosep]
  \item \textbf{Conductos} --- tuberías y canales abiertos de sección
        transversal arbitraria, descritos por la ecuación de Manning y un
        conjunto de tablas de geometría de sección transversal.
  \item \textbf{Bombas} --- gobernadas por curvas características (cinco tipos
        de curva o un modelo de bomba ideal) con histéresis de arranque/parada
        según el tirante.
  \item \textbf{Orificios} --- aberturas de fondo o laterales con
        comportamiento de descarga tipo vertedero/orificio.
  \item \textbf{Vertederos} --- transversales, laterales, de escotadura en V o
        trapezoidales, con transición opcional a escurrimiento por orificio
        bajo sumersión.
  \item \textbf{Descargas (outlets)} --- curvas de gasto en función de la
        carga o del tirante (tabulares o funcionales).
\end{itemize}

\subsection{Variables de estado y unidades internas}

En el tránsito de caudales las variables de estado fundamentales son
exactamente tres (Tabla~\ref{tab:statevars}):

\begin{table}[H]
\centering
\begin{tabular}{ll}
\toprule
\textbf{Símbolo} & \textbf{Descripción} \\
\midrule
$H$  & Carga hidráulica del agua en un nodo \\
$Q$  & Caudal en un enlace \\
$A$  & Área hidráulica en un enlace (inferida de $Q$ y de la geometría) \\
\bottomrule
\end{tabular}
\caption{Las tres variables de estado fundamentales del tránsito de caudales.}
\label{tab:statevars}
\end{table}

Todas las demás magnitudes se derivan de estas tres variables, de las entradas
externas y de los parámetros fijos. Internamente el motor realiza todos los
cálculos en \textbf{pies} (longitud), \textbf{segundos} (tiempo), \textbf{cfs}
(caudal), \textbf{ft$^2$} (área) y \textbf{ft$^3$} (volumen). Los valores de
entrada del proyecto en unidades SI (métricas) o US se convierten a estas
unidades internas al momento de la lectura; los factores de conversión son
\begin{align}
  \mathrm{Ucf}[\mathrm{LENGTH}] &=
    \begin{cases} 1.0 & \text{(US)} \\ 0.3048 & \text{(SI)} \end{cases}, \\
  \mathrm{Ucf}[\mathrm{VOLUME}] &=
    \begin{cases} 1.0 & \text{(US)} \\ 0.02832 & \text{(SI)} \end{cases},
\end{align}
(obsérvese que el factor SI de volumen es el valor \emph{truncado} del motor
legado, no $0.3048^3$, conservado deliberadamente para mantener la paridad de
punto flotante con el motor SWMM legado). La conversión de caudal usa
$\mathrm{Qcf} = \{1.0,\ 448.831,\ 0.64632,\ 0.02832,\ 28.317,\ 2.4466\}$ para
cfs/gpm/MGD/cms/LPS/MLD.

\subsection{Las cuatro formulaciones de tránsito}

El motor selecciona el método de tránsito de caudales con la opción
\texttt{FLOW\_ROUTING}. La Tabla~\ref{tab:routing} resume los cuatro métodos.

\begin{table}[H]
\centering
\begin{tabular}{@{}l p{4.6cm} p{5.4cm}@{}}
\toprule
\textbf{Método} & \textbf{Ecuaciones de gobierno} & \textbf{Características} \\
\midrule
Permanente & continuidad + tirante normal de Manning & pasante, sin almacenamiento/remanso \\
Onda cinemática & continuidad + curva de gasto de escurrimiento uniforme & sin remanso, inversión ni presión \\
Onda dinámica (defecto) & St.\ Venant completo & remanso, pérdidas, inversión, presión \\
Volúmenes finitos (explícito) & St.\ Venant, Godunov/HLLC & flujo transcrítico, seco/húmedo nativo \\
\bottomrule
\end{tabular}
\caption{Las cuatro formulaciones de tránsito 1D del motor.}
\label{tab:routing}
\end{table}

La onda dinámica es la opción por defecto y aquí se trata con mayor profundidad.
Cada método de tránsito se invoca desde un orquestador común \texttt{Router}.

% =============================================================================
% 2. COMPONENTES DEL MODELO HIDRÁULICO
% =============================================================================
\section{Componentes del modelo hidráulico}

\subsection{Datos de los nodos}

Cada nodo lleva propiedades geométricas estáticas y estado dinámico
(\texttt{NodeData.hpp}):
\begin{itemize}[nosep]
  \item $z_{\mathrm{inv}}$ --- cota de batea (ft).
  \item $d_{\mathrm{full}}$ --- profundidad total (tirante hasta el brocal /
        borde superior del estanque).
  \item $d_{\mathrm{sur}}$ --- profundidad máxima admisible de escurrimiento
        en carga sobre la profundidad total; el límite de anegamiento es
        $z_{\mathrm{inv}} + d_{\mathrm{full}} + d_{\mathrm{sur}}$.
  \item $A_{\mathrm{pond}}$ --- área de encharcamiento usada cuando el nodo se
        anega sobre el brocal y el encharcamiento está habilitado.
  \item $z_{\mathrm{crown}}$ --- cota de la clave del conducto más alto
        conectado; el umbral de entrada en carga es $y_{\mathrm{crown}} =
        z_{\mathrm{crown}} - z_{\mathrm{inv}}$.
  \item $\mathrm{deg}$ --- número de enlaces conectados; si es negativo, el
        nodo es un nodo terminal aguas arriba.
  \item $V_{\mathrm{full}}$ --- volumen a profundidad total.
\end{itemize}

Estado dinámico (actualizado en cada paso de tránsito):
\begin{itemize}[nosep]
  \item $d$ --- tirante de agua sobre la batea (ft);
  \item $H = z_{\mathrm{inv}} + d$ --- carga hidráulica (ft);
  \item $V$ --- volumen almacenado (ft$^3$);
  \item $Q_{\mathrm{lat}}$ --- aporte lateral total (ft$^3$/s);
  \item $Q_{\mathrm{in}}$, $Q_{\mathrm{out}}$ --- ingreso y salida totales;
  \item $Q_{\mathrm{ovfl}}$ --- tasa de rebalse/anegamiento;
  \item $Q_{\mathrm{loss}}$ --- tasa de pérdidas por evaporación +
        filtración/exfiltración;
  \item $\Delta Q_{\mathrm{net}}^{t-1}$ --- ingreso neto $Q_{\mathrm{in}} -
        Q_{\mathrm{out}}$ del paso anterior (usado para el promediado
        trapezoidal);
  \item valores del paso anterior $d^{t-1}$, $V^{t-1}$, $Q_{\mathrm{lat}}^{t-1}$.
\end{itemize}

\subsubsection{Funciones de área superficial y volumen de los nodos}

Las funciones de geometría de nodo están implementadas en
\src{hydraulics/Node.cpp}.

\textbf{Volumen de cámara de unión / emisario / divisor}: con
$\mathrm{MIN\_SURFAREA} = 12.566\ \sfft$ ($\approx 4\pi$, el área de una
cámara de inspección de 4 pies),
\begin{equation}
  V(d) = V_{\mathrm{full}}\,\frac{d}{d_{\mathrm{full}}},
\end{equation}
donde $V_{\mathrm{full}} = \mathrm{MIN\_SURFAREA}\cdot d_{\mathrm{full}}$
salvo que se haya redefinido.

\textbf{Volumen del estanque de almacenamiento}:
\begin{itemize}[nosep]
  \item Tabular: $V = \mathrm{table}(d\cdot \mathrm{Ucf}[L])/\mathrm{Ucf}[V]$
        con interpolación de tabla;
  \item Las formas geométricas (cilíndrica, cónica, paraboloide, piramidal)
        usan la relación cuadrática de área superficial
        $A(d) = c + a\,d + b\,d^{2}$, integrada al polinomio cúbico
        \begin{equation}
          V(d) = d\left(c + d\!\left(\tfrac{a}{2} + d\,\tfrac{b}{3}\right)\right);
        \end{equation}
  \item Almacenamiento funcional: ley de potencia $A(d) = c + a\,d^{b}$
        integrada analíticamente.
\end{itemize}

\textbf{Área superficial} $A_{\mathrm{surf}}(d)$: los nodos no-estanque
retornan 0 (el piso de área mínima se aplica más adelante en la actualización
de nodo de la onda dinámica); los estanques retornan el valor de la tabla
(extrapolando linealmente) o la forma analítica $A(d) = c + a\,d + b\,d^{2}$ /
$A(d) = c + a\,d^{b}$.

\textbf{Área de encharcamiento}:
\begin{equation}
  A_{\mathrm{pond}}(d) =
  \begin{cases}
    A_{\mathrm{surf}}(d) & d \le d_{\mathrm{full}} \text{ o } A_{\mathrm{pond}} = 0,\\
    A_{\mathrm{pond}} & d > d_{\mathrm{full}} \text{ (anegado)}.
  \end{cases}
\end{equation}

\textbf{Tirante a partir del volumen} (la inversa de la función de volumen,
usada por la actualización de estanque de la onda cinemática y por los
reportes): los nodos cámara de unión usan $d = V/\mathrm{MIN\_SURFAREA}$; los
estanques invierten las relaciones tabular/geométrica/funcional, usando
Newton--Raphson con fallback de bisección cuando no existe forma cerrada.

\textbf{Caudal máximo de salida} (usado por la reunión de ingreso de enlace de
la onda cinemática y del escurrimiento permanente):
\begin{equation}
  Q_{\max} = Q_{\mathrm{in}} + \frac{V_{\mathrm{old}}}{dt},
  \qquad Q_{\mathrm{in}} = \min\!\left(Q_{\mathrm{in}}, Q_{\max}\right).
\end{equation}

\subsection{Datos de los enlaces}

Cada enlace lleva propiedades estáticas y estado dinámico
(\texttt{LinkData.hpp}). Para los conductos, los parámetros de capacidad de
conducción por enlace (por barril) son (\texttt{Link.cpp},
\texttt{computeConveyance}):
\begin{align}
  \beta &= \frac{\mathrm{PHI}\,\sqrt{|S_0|}}{n}, \qquad \mathrm{PHI} = 1.486
    \quad\text{(factor US de Manning)}, \\
  \mathrm{rough\_factor} &= g\left(\frac{n}{\mathrm{PHI}}\right)^{2},\\
  Q_{\mathrm{full}} &= S_{\mathrm{full}}\cdot\beta,
  \qquad Q_{\max} = S_{\max}\cdot\beta,
\end{align}
donde $n$ es la rugosidad de Manning, $S_0$ la pendiente del conducto,
$S(a) = A\,R^{2/3}$ el \emph{factor de sección} (ver Sección~\ref{sec:xsect}),
$S_{\mathrm{full}} = S(A_{\mathrm{full}})$ y $S_{\max} = \max_a S(a)$.

Estado dinámico por enlace: $Q$ (caudal actual, $+$ = nodo1$\to$nodo2), $d$
(tirante medio), $V$ (volumen), $\mathrm{Fr}$ (número de Froude),
$\mathrm{flow\_class}$ (DRY, UP\_DRY, DN\_DRY, SUBCRITICAL, SUPERCRITICAL,
UP\_CRITICAL, DN\_CRITICAL) y valores del paso anterior.

\subsection{Pendiente y longitud del conducto}

Para un conducto que conecta los nodos $i$ y $j$ con elevaciones de extremo
$z_1, z_2$, las cotas de batea de los extremos del conducto son $Z_1 =
z_{\mathrm{inv},1} + z_1$, $Z_2 = z_{\mathrm{inv},2} + z_2$. La pendiente del
conducto es
\begin{equation}
  S_0 = \frac{\Delta y}{\Delta x},
  \qquad \Delta y = Z_1 - Z_2, \quad
  \Delta x = \sqrt{L^{2} - \Delta y^{2}},
\end{equation}
con un desnivel mínimo impuesto de $\lvert \Delta y \rvert \ge 0.001\ \ft$
(sobrescribible).

\subsection{Alargamiento de conductos (estabilidad de Courant)}

Si se entrega un \texttt{LENGTHENING\_STEP} positivo, los conductos cortos se
alargan artificialmente para que se cumpla la condición de Courant para el paso
de tiempo entregado por el usuario (\texttt{Routing.cpp:135--195}):
\begin{equation}
  t_{\mathrm{step}} = \min(\mathrm{routing\_step},\ \mathrm{lengthening\_step}),
  \qquad
  v_{\mathrm{full}} = \frac{\mathrm{PHI}}{n}\,\frac{S_{\mathrm{full}}}
    {\sqrt{|S_0|}\,A_{\mathrm{full}}},
\end{equation}
\begin{equation}
  \mathrm{ratio} = \frac{\left(\sqrt{g\,y_{\mathrm{full}}} + v_{\mathrm{full}}\right)
    t_{\mathrm{step}}}{L},
  \qquad
  L' = \max(L,\ \mathrm{ratio}\cdot L),
\end{equation}
En canales abiertos, $y_{\mathrm{full}}$ se reemplaza por el tirante hidráulico
$A_{\mathrm{full}}/W_{\max}$. Si $L' > L$, la pendiente y la rugosidad se
ajustan para preservar igual pérdida de carga a cualquier caudal,
\begin{equation}
  S_0' = \frac{|S_0|}{L'/L}, \qquad
  n' = \frac{n}{\sqrt{L'/L}},
\end{equation}
y $\beta$, $\mathrm{rough\_factor}$, $Q_{\mathrm{full}}$, $Q_{\max}$ se
recalculan a partir de los valores ajustados. La longitud alargada $L'$
(\texttt{mod\_length}) se usa en la ecuación de momentum; la longitud cruda $L$
se usa para la contabilidad de volúmenes.

% =============================================================================
% 3. ORQUESTACIÓN DE LA SIMULACIÓN
% =============================================================================
\section{Orquestación de la simulación}

\subsection{El ciclo de avance temporal global}

El motor es una API de paso único. Cada llamada a \texttt{SWMMEngine::step()}
ejecuta exactamente un paso de tránsito (o un paso solo de escorrentía cuando
el tránsito está deshabilitado). La secuencia dentro de un paso es
(\texttt{SWMMEngine.cpp:941}):
\begin{enumerate}[nosep]
  \item Calcular el paso de tránsito $dt_{\mathrm{next}}$. Si la opción de paso
        variable está habilitada, el paso limitado por Courant es
        \[
          dt_{\mathrm{cfl}} = \texttt{router.getAdaptiveStep}(
            \texttt{routing\_step},\ \texttt{variable\_step}),
        \]
        luego $dt_{\mathrm{next}} = \min(\texttt{routing\_step},\ dt_{\mathrm{cfl}})$
        acotado a la duración restante de la simulación.
  \item Guardar el estado hidráulico anterior (\texttt{ctx.save\_state()}).
  \item Reiniciar los acumuladores de balance de masa del paso.
  \item Ejecutar la cadena:
        \begin{enumerate}[nosep]
          \item \texttt{stepRunoff($dt$)} --- hidrología (siempre);
          \item si el tránsito está habilitado: \texttt{step\-Routing($dt$)}
                (hidráulica + calidad), \texttt{update\-Statistics($dt$)},
                \texttt{update\-Routing\-Mass\-Balance($dt$)};
          \item \texttt{compute\-Final\-Storage()}.
        \end{enumerate}
  \item Avanzar el reloj de la simulación y publicar las instantáneas de
        salida.
\end{enumerate}

\subsection{Relojes de escorrentía y de tránsito}

La hidrología y la hidráulica corren en relojes \emph{independientes}:
\begin{itemize}[nosep]
  \item El \textbf{reloj de escorrentía} avanza al paso húmedo (300\,s por
        defecto) o al paso seco (3600\,s por defecto), acortado para alinear
        con los límites de los pluviómetros. \texttt{stepRunoff} ejecuta
        múltiples subpasos de escorrentía dentro de un paso de tránsito
        ($\text{mientras } t_{\mathrm{runoff}} < t_{\mathrm{routing}} + dt$).
  \item El \textbf{reloj de tránsito} avanza en $dt$ (el paso de tránsito) en
        cada llamada a \texttt{step()}.
\end{itemize}
Después de los subpasos de escorrentía, el aporte lateral de tiempo húmedo
entregado a los nodos se obtiene por \emph{interpolación lineal} entre las
evaluaciones de escorrentía que lo enmarcan: con $f = (t_{\mathrm{elapsed}} -
t_{\mathrm{runoff,old}}) / (t_{\mathrm{runoff,new}} - t_{\mathrm{runoff,old}})$,
\begin{equation}
  Q_{\mathrm{runoff}} = (1-f)\,Q_{\mathrm{runoff}}^{\mathrm{old}} +
  f\,Q_{\mathrm{runoff}}^{\mathrm{new}},
\end{equation}
más los aportes por escurrimiento superficial recibido (runon) y por napa
freática. Estos se distribuyen en \texttt{nodes.runoff\_inflow} /
\texttt{nodes.gw\_inflow}.

\subsection{Aportes laterales}

En cada paso de tránsito, los aportes laterales se ensamblan a partir de los
búferes de fuente descompuestos en el orden legado
(\texttt{assembleLateralInflows}, \texttt{SWMMEngine.cpp:5691}):
\begin{equation}
  Q_{\mathrm{lat}} = Q_{\mathrm{ext}} + Q_{\mathrm{dwf}} +
  Q_{\mathrm{wet}} + Q_{\mathrm{gw}} + Q_{\mathrm{rdii}} +
  Q_{\mathrm{iface}} + Q_{\mathrm{user}} + Q_{\mathrm{coupling}},
\end{equation}
donde los términos son, respectivamente: series de tiempo externas/ingresos,
caudal de tiempo seco (DWF), ingreso de tiempo húmedo (escorrentía), ingreso de
napa freática, RDII (infiltración/entrada por lluvia), ingreso por archivo de
interfaz, ingreso forzado por el usuario e ingreso por acoplamiento 1D--2D.

\subsection{El paso de tránsito}

\texttt{Router::step(ctx, dt, evap\_rate, non\_conduit\_fn)}
(\texttt{Routing.cpp:308}) ejecuta, en orden:
\begin{enumerate}[nosep]
  \item \textbf{initNodeFlows} --- inicializa el ingreso de cada nodo desde el
        aporte lateral y la salida desde las pérdidas (evaporación del estanque
        y exfiltración de Green--Ampt, acotadas conjuntamente por el volumen
        almacenado); fija el rebalse a partir del exceso de volumen almacenado.
  \item \textbf{computeConduitLosses} --- tasas de pérdida de los conductos por
        evaporación (secciones abiertas) y filtración. Para DYNWAVE esto se
        recalcula en su lugar en cada iteración de Picard dentro del solver.
  \item \textbf{setAllOutfallDepths} --- fija los niveles de agua de borde de
        los emisarios (Sección~\ref{sec:outfall}).
  \item \textbf{Despacho del solver} --- KINWAVE, DYNWAVE, FV o STEADY.
  \item \textbf{computeDividerFlows} --- aplica la lógica de divisores de
        flujo.
  \item \textbf{updateLinkStates} --- recalcula las cargas de los nodos
        $H = z_{\mathrm{inv}} + d$.
\end{enumerate}

Dentro del solver de onda dinámica, la función de retorno de no-conducto
calcula los caudales de bombas/orificios/vertederos/descargas dentro de la
iteración de Picard, en orden de índice de enlace, con distribución inmediata
en los acumuladores de ingreso/salida de los nodos.

% =============================================================================
% 4. GEOMETRÍA DE LA SECCIÓN TRANSVERSAL
% =============================================================================
\section{Geometría de la sección transversal}
\label{sec:xsect}

\subsection{Las cuatro relaciones geométricas}

Para un conducto, el área hidráulica $A$, el espejo de agua (ancho
superficial) $W$, el perímetro mojado $P_w$ y el radio hidráulico $R = A/P_w$
son funciones del tirante $y$. El motor precalcula, para cada forma, las cuatro
relaciones fundamentales (\texttt{XSectKernels.hpp},
\texttt{xsect\_tables.hpp}):
\begin{align}
  A &= A(y), \qquad W = W(y), \qquad R = R(y), \\
  S &= S(A) = A\,R^{2/3} \quad \text{(factor de sección, ``forma de
  capacidad de conducción''),}
\end{align}
junto con sus inversas $y = y(A)$, $A = A(S)$ y la derivada $dS/dA$. Existen
expresiones cerradas para las formas estándar (circular, circular llena,
rectangular, trapezoidal, triangular, parabólica, potencia, huevo, herradura,
gótica, canasta, etc.); las formas irregulares y personalizadas usan tablas
(transectas) con interpolación lineal, y la inversa $A(S)$ se resuelve por
iteración de Newton o inversión de tabla.

\subsection{Manning y el factor de sección}

La ecuación de Manning en forma US es
\begin{equation}
  Q = \frac{1.486}{n}\,A\,R^{2/3}\,\sqrt{S_f}
  = \beta\,\Psi(A)\,\sqrt{S_f/S_0},
  \qquad \Psi(A) = A\,R^{2/3},
\end{equation}
donde $S_f$ es la pendiente de fricción, $\beta = (1.486/n)\sqrt{S_0}$ y
$\Psi(A)$ es el factor de sección. Para muchas formas cerradas el factor de
sección presenta un máximo bajo la profundidad total, de modo que el caudal
máximo ocurre a un área menor que la de sección llena.

% =============================================================================
% 5. TRÁNSITO DE ESCURRIMIENTO PERMANENTE
% =============================================================================
\section{Tránsito de escurrimiento permanente}

Bajo escurrimiento permanente (\texttt{executeSteadyFlow},
\texttt{Routing.cpp:637}), los enlaces se procesan en orden topológicamente
ordenado. Para cada conducto:
\begin{enumerate}[nosep]
  \item Se reúne el ingreso del nodo aguas arriba $Q_{\mathrm{in}}$ (limitado
        por el caudal máximo de salida).
  \item El caudal por barril es $q = Q_{\mathrm{in}}/\mathrm{barrels}$ menos la
        tasa de pérdida del conducto, acotado a cero.
  \item Si $q \ge Q_{\mathrm{full}}$: $q = Q_{\mathrm{full}}$ y
        $A = A_{\mathrm{full}}$; en caso contrario, el área de tirante normal
        se obtiene del factor de sección inverso,
        \begin{equation}
          s = \frac{q}{\beta}, \qquad A = A(s),
        \end{equation}
        y el tirante de $y = y(A)$.
  \item El escurrimiento es uniforme a lo largo del conducto ($Q_{\mathrm{out}}
        = q\cdot\mathrm{barrels}$, $A_1 = A_2 = A$); el volumen del enlace es
        $V = A\,L\,\mathrm{barrels}$.
\end{enumerate}
Los enlaces no-conducto simplemente dejan pasar el ingreso del nodo aguas
arriba. El escurrimiento permanente converge en una sola pasada.

% =============================================================================
% 6. ONDA CINEMÁTICA
% =============================================================================
\section{Tránsito por onda cinemática}

\subsection{Ecuaciones de gobierno}

El modelo de onda cinemática combina la ecuación de continuidad con la curva
de gasto de escurrimiento uniforme, eliminando los términos de inercia,
gradiente de presión y (parcialmente) convectivos de la ecuación de momentum de
St.\ Venant. No puede representar remanso, inversión del flujo ni
presurización, y requiere una red dirigida acíclica. Partiendo de las
ecuaciones de St.\ Venant y fijando la pendiente de fondo igual a la pendiente
de fricción, $S_0 = S_f$, el par de gobierno es:
\begin{equation}
  \frac{\partial A}{\partial t} + \frac{\partial Q}{\partial x} = 0
  \qquad\text{(continuidad),}
\end{equation}
\begin{equation}
  Q = \beta\,\Psi(A), \qquad \Psi(A) = A\,R^{2/3},
  \qquad \beta = \frac{1.486}{n}\sqrt{S_0}
  \qquad\text{(curva de gasto).}
\end{equation}

\subsection{Esquema de diferencias finitas}

Un esquema de diferencias finitas implícito de Wendroff ponderado discretiza la
ecuación de continuidad entre los extremos aguas arriba (1) y aguas abajo (2):
\begin{equation}
  \frac{(1-\theta)(A_1^{t+1}-A_1^{t}) +
        \theta\,(A_2^{t+1}-A_2^{t})}{\Delta t}
  \;+\;
  \frac{(1-\varphi)(Q_2^{t}-Q_1^{t}) +
        \varphi\,(Q_2^{t+1}-Q_1^{t+1})}{L}
  = 0,
\end{equation}
con $\theta = \varphi = 0.6$. Como cada cámara de unión tiene a lo más un
conducto de salida, procesar los enlaces en orden topológico deja solo a
$A_2^{t+1}$ y $Q_2^{t+1}$ como incógnitas; $Q_1^{t+1}$ se conoce desde el
ingreso del nodo aguas arriba y $A_1^{t+1}$ desde el factor de sección inverso
en $Q_1^{t+1}/\beta$. Sustituyendo la curva de gasto se obtiene la única
ecuación no lineal
\begin{equation}
  f\!\left(A_2^{t+1}\right) = \beta\,\Psi\!\left(A_2^{t+1}\right) +
  C_1\,A_2^{t+1} + C_2 = 0,
\end{equation}
con (en la implementación normalizada del motor,
\texttt{KinematicWave.cpp:solveConduit})
\begin{align}
  C_1 &= \frac{L\,\theta}{\Delta t\,\varphi},\\
  C_2 &= \frac{L}{\Delta t\,\varphi}
        \Big[(1-\theta)(A_1^{t+1} - A_1^t) - \theta\,A_2^t\Big]
        + \frac{1-\varphi}{\varphi}(Q_2^t - Q_1^t) - Q_1^{t+1}.
\end{align}

\subsection{Búsqueda de la raíz}

La ecuación $f(A_2^{t+1}) = 0$ se resuelve con una iteración de Newton--Raphson
enmarcada (máximo 40 iteraciones). Como el factor de sección puede tener dos
raíces (máximo en $A_{\max}$ bajo sección llena), el marco se preselecciona:
el marco inicial es $[A_{\max}, A_{\mathrm{full}}]$; si ambos límites producen
el mismo signo se reinicia a $[0, A_{\max}]$. Si ambos límites dan $f$ negativo
el conducto está a sección llena ($Q = Q_{\mathrm{full}}$); si ambos dan $f$
positivo el caudal es cero.

\subsection{Estanques bajo onda cinemática}

Los estanques de almacenamiento pueden tener cualquier número de enlaces de
salida. Su balance de masa se integra con la regla trapezoidal,
\begin{equation}
  V^{t+1} = V^{t} + \tfrac{1}{2}\left(Q_{\mathrm{in}}^{t} + Q_{\mathrm{in}}^{t+1}
  - Q_{\mathrm{out}}^{t} - Q_{\mathrm{out}}^{t+1}\right)\Delta t,
\end{equation}
resuelto por aproximaciones sucesivas ($\omega = 0.55$, tolerancia de
convergencia 0.005 ft) porque tanto $Q_{\mathrm{out}}$ como $V$ dependen de la
carga. La carga del estanque se actualiza una vez más después de que se
conocen todos los caudales de los enlaces (\texttt{Routing.cpp:345--389}).

% =============================================================================
% 7. ONDA DINÁMICA: ECUACIONES DE GOBIERNO
% =============================================================================
\section{Onda dinámica: ecuaciones de gobierno}

\subsection{Las ecuaciones de St.\ Venant}

El modelo de onda dinámica resuelve las ecuaciones unidimensionales completas
de de~Saint-Venant (continuidad y momentum):
\begin{equation}
  \underbrace{\frac{\partial A}{\partial t}}_{\text{almacenamiento}} +
  \underbrace{\frac{\partial Q}{\partial x}}_{\text{flujo convectivo}}
  = 0,
  \label{eq:sv-cont}
\end{equation}
\begin{equation}
  \underbrace{\frac{\partial Q}{\partial t}}_{\text{aceleración local}} +
  \underbrace{\frac{\partial}{\partial x}\!\left(\frac{Q^{2}}{A}\right)}_{\text{aceleración convectiva}}
  + \underbrace{gA\,\frac{\partial H}{\partial x}}_{\text{gradiente de presión}}
  + \underbrace{gA\,S_f}_{\text{fricción}}
  = 0,
  \label{eq:sv-mom}
\end{equation}
donde $A$ es el área hidráulica, $Q$ el caudal, $H = Z + y$ la carga hidráulica
($Z$ cota de batea del conducto, $y$ tirante), $g$ la gravedad y $S_f$ la
pendiente de fricción. La pendiente de fricción sigue a Manning:
\begin{equation}
  S_f = \left(\frac{n}{1.486}\right)^{2}
  \frac{Q\,\lvert U \rvert}{A\,R^{4/3}},
  \qquad U = \frac{Q}{A}.
\end{equation}

\subsection{Ecuación de momentum en diferencias finitas}

El motor usa una forma implícita (Euler hacia atrás) en diferencias finitas.
Con diferencias espaciales sobre la longitud del conducto $L$ y diferencias
temporales sobre $\Delta t$, la actualización de caudal en cada conducto es
\begin{equation}
  Q^{t+\Delta t}
  = \frac{Q^{t} + \Delta Q_{\text{inercia}} + \Delta Q_{\text{presión}}
        + \Delta Q_{\text{pérdida}}}
         {1 + \Delta Q_{\text{fricción}}},
  \label{eq:mom-update}
\end{equation}
donde todas las magnitudes geométricas se evalúan en el nuevo tiempo
$t + \Delta t$. Los cuatro términos son:
\begin{align}
  \Delta Q_{\text{inercia}}
    &= 2\bar{U}\left(\bar{A}^{t+\Delta t} - \bar{A}^{t}\right)
       + \bar{U}^{2}\,\frac{A_2 - A_1}{L}\,\Delta t, \label{eq:mom-inertia}\\
  \Delta Q_{\text{presión}}
    &= -g\bar{A}\,\frac{H_2 - H_1}{L}\,\Delta t, \label{eq:mom-pressure}\\
  \Delta Q_{\text{fricción}}
    &= g\left(\frac{n}{1.486}\right)^{2}
       \frac{\lvert \bar{U} \rvert}{\bar{R}^{4/3}}\,\Delta t, \label{eq:mom-friction}\\
  \Delta Q_{\text{pérdida}}
    &= \text{pérdidas locales (menores) y términos de evaporación/filtración}.
\end{align}

\subsection{Continuidad en los nodos}

La conservación de volumen en cada ensamble de nodo (el nodo más la mitad de
cada enlace conectado) requiere
\begin{equation}
  \frac{\partial V}{\partial t}
  = A_{S}\,\frac{\partial H}{\partial t} = \sum Q,
\end{equation}
con $A_S = A_{SN} + \sum A_{SL}$ el área superficial del ensamble (el área
superficial propia del nodo más las áreas superficiales de medio enlace que
aporta cada conducto). En forma de diferencias finitas:
\begin{equation}
  H^{t+\Delta t} = H^{t} +
  \frac{\tfrac{\Delta t}{2}\left(\sum Q^{t} + \sum Q^{t+\Delta t}\right)}
       {\left(A_{SN} + \sum A_{SL}\right)^{t+\Delta t}},
\end{equation}
es decir, una actualización trapezoidal de la carga usando el \emph{promedio}
de los caudales netos en $t$ y $t+\Delta t$.

% =============================================================================
% 8. ALGORITMO DE SOLUCIÓN DE LA ONDA DINÁMICA
% =============================================================================
\section{Algoritmo de solución de la onda dinámica}

\subsection{Resumen: la iteración de Picard (aproximaciones sucesivas)}

Las ecuaciones~\eqref{eq:mom-update} y la continuidad de nodo se resuelven
implícitamente sobre cada paso de tiempo mediante \emph{iteración funcional}
(iteración de Picard / aproximaciones sucesivas). Los pasos, que coinciden con
el \texttt{DWSolver::execute} del motor (\texttt{DynamicWave.cpp:1026}), son:

\begin{enumerate}[nosep]
  \item \textbf{initNodeStates} --- reinicia los acumuladores por nodo: ingreso
        = 0, salida = pérdidas $+$ aporte lateral negativo, ingreso $+$ aporte
        lateral positivo; área superficial desde la geometría de
        almacenamiento/encharcamiento; convergido = 0, $\sum dQ/dH = 0$.
  \item \textbf{computeLinkGeometry} --- calcula en lote los tirantes de
        extremo $y_1, y_2$, el tirante medio $\bar y$, los espejos de agua, las
        áreas y los radios hidráulicos de todos los conductos desde las cargas
        actuales de los nodos; clasifica el régimen de escurrimiento de cada
        conducto; aplica las correcciones de ranura de Preissmann donde hay
        escurrimiento en carga.
  \item \textbf{Núcleos de momentum} --- resuelve la ecuación de momentum
        implícita para cada conducto (Sección~\ref{sec:momentum-kernel}),
        registra los nuevos caudales y los distribuye en los acumuladores de
        ingreso/salida de los nodos, acumulando $\sum dQ/dH$ y las
        contribuciones de área superficial.
  \item \textbf{Estructuras no-conducto} --- calcula los caudales de
        bombas/\allowbreak orificios/\allowbreak vertederos/\allowbreak
        descargas (Sección~\ref{sec:structures}).
  \item \textbf{Niveles de emisario} --- vuelve a fijar los niveles de borde de
        los emisarios desde los caudales actuales.
  \item \textbf{Actualización de tirante en los nodos} --- resuelve la ecuación
        de continuidad para la carga de cada nodo
        (Sección~\ref{sec:node-continuity}); cuenta los nodos no convergidos.
  \item Si algún nodo cambió más que la tolerancia de carga y no se ha alcanzado
        el límite de intentos, marca como \emph{omitidos (bypassed)} los
        enlaces cuyos dos nodos extremos convergieron (sus caudales se
        mantienen), e itera.
\end{enumerate}

Subrelajación: después de la primera iteración, todo caudal y carga
actualizados se relajan con $\omega = 0.5$,
\begin{equation}
  x^{k+1} = (1-\omega)\,x^{k} + \omega\,x^{k+1}_{\mathrm{crudo}},
  \qquad \omega = 0.5,
\end{equation}
y un cambio de signo en el caudal se fuerza a pasar por un valor pequeño no
nulo ($q = \pm 0.001$).

\subsection{Convergencia}

Un nodo está convergido cuando tanto el residual de Picard crudo como el
movimiento aceptado (posiblemente mezclado por Anderson) están dentro de la
tolerancia de carga $\varepsilon_H = 0.005\ \ft$ (por defecto). El paso de
tránsito converge cuando ningún nodo no-emisario queda sin converger. El número
máximo de intentos es 8 por defecto (\texttt{MAX\_TRIALS}); un paso que
converge en su último intento permitido aún se cuenta como convergido.

\subsection{Clasificación del escurrimiento}

Durante el cálculo de la geometría, cada conducto se clasifica en una de siete
clases de escurrimiento (DRY, UP\_DRY, DN\_DRY, SUBCRITICAL, SUPERCRITICAL,
UP\_CRITICAL, DN\_CRITICAL) en base a los tirantes de extremo y a la carga
relativa a las cotas de batea de los extremos, usando el tirante normal $y_n$
(desde el factor de sección inverso en $Q/\beta$) y el tirante crítico $y_c$
(Sección~\ref{sec:outfall}). La clase determina:
\begin{itemize}[nosep]
  \item los ajustes de área superficial para extremos secos/críticos
        (Tabla~\ref{tab:surfadj}),
  \item el tirante usado en el extremo crítico (reemplazado por
        $\min(y_n, y_c)$),
  \item el límite de caudal normal basado en la fricción.
\end{itemize}

\begin{table}[H]
\centering
\small
\begin{tabular}{lll}
\toprule
\textbf{Condición} & \textbf{Criterio} & \textbf{Ajuste} \\
\midrule
Seco aguas arriba & $y_1 = 0$; $Z_1 > E_1$ & $A_{SL1} = 0$ si $H_2 \le Z_1$, si no ajuste crítico \\
Seco aguas abajo & $y_2 = 0$; $Z_2 > E_2$ & $A_{SL2} = 0$ si $H_1 \le Z_2$, si no ajuste crítico \\
Crítico aguas arriba & $Q < 0$; $Z_1 > E_1$; $H_1 - Z_1 < y_*$ & $y_1 = y_*$; $A_{SL1} = 0$ \\
Crítico aguas abajo & $Q > 0$; $Z_2 > E_2$; $H_2 - Z_2 < y_*$ & $y_2 = y_*$; $A_{SL2} = 0$ \\
\bottomrule
\end{tabular}
\caption{Ajustes de área superficial y de tirante para extremos de conducto
secos/críticos ($y_* = \min(y_n, y_c)$; $E$ = batea del nodo, $Z$ = batea del
conducto).}
\label{tab:surfadj}
\end{table}

\subsection{Contribuciones de área superficial}

Para un conducto subcrítico, las áreas de superficie libre aportadas a sus dos
nodos extremos son
\begin{equation}
  A_{SL1} = \frac{W_1 + \bar{W}}{2}\,\frac{L}{2},
  \qquad
  A_{SL2} = \frac{\bar{W} + W_2}{2}\,\frac{L}{2}\,
  \mathrm{fasnh},
\end{equation}
donde $W_1, W_2, \bar W$ son los espejos de agua en $y_1, y_2, \bar y$, y
$\mathrm{fasnh}$ es un factor entre el tirante normal y el crítico para un
extremo aguas abajo casi crítico. Para extremos secos/críticos el área
superficial de medio enlace se ajusta según la Tabla~\ref{tab:surfadj}; un
conducto completamente seco aporta $A_{SL1} = A_{SL2} =
\mathrm{FUDGE}\cdot L/2$. Un conducto cerrado cuyo tirante medio supera el
corte de clave ($y/y_{\mathrm{full}} \ge 0.96$ bajo EXTRAN) usa un ancho
acotado en la clave en el cálculo del área superficial, entregando la pequeña
contribución no nula que mantiene acotado el denominador de Picard del nodo.

% =============================================================================
% 9. EL NÚCLEO DE MOMENTUM DEL ENLACE
% =============================================================================
\section{El núcleo de momentum del enlace}
\label{sec:momentum-kernel}

Cada conducto se resuelve con uno de los núcleos de momentum de
\texttt{DynamicWave.cpp} (\texttt{process\-Manning\-Link},
\texttt{process\-Dry\-Link}, \texttt{process\-Force\-Main\-Link}), despachado
por categoría de momentum. Esta sección transcribe el núcleo de Manning
(superficie libre y conducto cerrado), que es el principal.

\subsection{Velocidad, número de Froude y amortiguación inercial}

Con $\bar A$ el área media y $q_{\mathrm{last}}$ el caudal por barril de la
iteración anterior,
\begin{equation}
  v = \frac{q_{\mathrm{last}}}{\bar A}, \qquad
  \lvert v \rvert \le V_{\max} = 50\ \ft/\st
  \quad\text{(tope de velocidad).}
\end{equation}
Para un conducto sin escurrimiento en carga, el número de Froude usa el tirante
hidráulico $\bar A/\bar W$:
\begin{equation}
  \mathrm{Fr} = \frac{\lvert v \rvert}{\sqrt{g\,\bar A / \bar W}}
  \quad\text{(0 para un conducto cerrado a menos de FUDGE de sección llena).}
\end{equation}
El factor de amortiguación inercial sigue el enfoque de \emph{inercia parcial
local} (mezcla lineal sobre $0.5 \le \mathrm{Fr} \le 1$):
\begin{equation}
  \sigma =
  \begin{cases}
    1 & \mathrm{Fr} \le 0.5,\\
    2\,(1 - \mathrm{Fr}) & 0.5 < \mathrm{Fr} < 1,\\
    0 & \mathrm{Fr} \ge 1,
  \end{cases}
  \qquad\text{es decir}\quad
  \sigma = \mathrm{clamp}\!\big(2(1-\mathrm{Fr}),\ 0,\ 1\big).
\end{equation}
La opción \texttt{INERTIAL\_DAMPING} sobrescribe esto: NONE fuerza $\sigma = 1$,
FULL fuerza $\sigma = 0$ (sin términos inerciales en absoluto); un conducto
cerrado con escurrimiento en carga siempre tiene $\sigma = 0$.

\subsection{Ponderación aguas arriba}

Los términos de presión y fricción usan área y radio hidráulico promedio
\emph{ponderados aguas arriba}, reflejando que el escurrimiento supercrítico
solo se ve influido por las condiciones aguas arriba. Con el factor de
ponderación
\begin{equation}
  \rho = \begin{cases} \sigma & q_{\mathrm{last}} > 0 \text{ y } H_1 \ge H_2
    \text{ (y no a sección llena),}\\ 1 & \text{en caso contrario,}\end{cases}
\end{equation}
\begin{equation}
  \bar A' = A_1 + \rho\left(\bar A - A_1\right),
  \qquad
  \bar R' = R_1 + \rho\left(\bar R - R_1\right).
\end{equation}

\subsection{Los seis términos de momentum}

El motor evalúa la ecuación de momentum como las seis contribuciones
$dq_1 \ldots dq_6$ (por barril):
\begin{align}
  dq_1 &= \Delta t \cdot \mathrm{rough\_factor} \cdot
        \frac{\lvert v \rvert}{\bar R'^{\,4/3}}, \label{eq:dq1}\\
  dq_2 &= \Delta t\, g\, \bar A'\, \frac{H_2 - H_1}{L}, \label{eq:dq2}\\
  dq_3 &= 2\,v\,\sigma\left(\bar A^{t+\Delta t} - \bar A^{t}\right), \label{eq:dq3}\\
  dq_4 &= \Delta t\, v^{2}\, \sigma\,\frac{A_2 - A_1}{L}, \label{eq:dq4}\\
  dq_5 &= \frac{\Delta t}{2L}
          \left[ K_{\mathrm{in}}\frac{\lvert q \rvert}{A_1} +
                K_{\mathrm{out}}\frac{\lvert q \rvert}{A_2} +
                K_{\mathrm{avg}}\frac{\lvert q \rvert}{\bar A} \right], \label{eq:dq5}\\
  dq_6 &= \frac{2.5\,\Delta t\, v}{L}\left(Q_{\mathrm{evap}} +
          Q_{\mathrm{filt}}\right), \label{eq:dq6}
\end{align}
donde $\mathrm{rough\_factor} = g(n/1.486)^2$, los coeficientes de pérdida
local $K_{\mathrm{in}}, K_{\mathrm{out}}, K_{\mathrm{avg}}$ son los
coeficientes de pérdida menor de entrada, salida y promedio (ponderado por
pérdidas) del usuario, y $Q_{\mathrm{evap}} + Q_{\mathrm{filt}}$ es la tasa de
pérdida del conducto.

La actualización de caudal (por barril) y su gradiente de carga son
\begin{equation}
  q = \frac{q_{\mathrm{old}} - dq_2 + dq_3 + dq_4 + dq_6}
           {1 + dq_1 + dq_5},
  \label{eq:q-update}
\end{equation}
\begin{equation}
  \frac{dQ}{dH} = \frac{1}{1 + dq_1 + dq_5}\,
  \frac{g\,\Delta t\, \bar A'}{L}\,\mathrm{barrels}.
  \label{eq:dqdh}
\end{equation}
El gradiente $dQ/dH$ alimenta el solver de nodos en carga
(Sección~\ref{sec:node-cont}).

\subsection{Limitación del caudal y postproceso}

Después de la actualización de momentum cruda, \texttt{applyFlowLimits}
(\texttt{DynamicWave.cpp:2211}) aplica, en orden:
\begin{enumerate}[nosep]
  \item \textbf{Control de entrada de alcantarilla} (FHWA HEC-5): si hay un
        código de alcantarilla y el conducto no está a sección llena,
        $q \leftarrow \min(q, q_{\mathrm{inlet}})$.
  \item \textbf{Límite de caudal normal}: para un conducto abierto/de
        superficie libre no lleno, si se cumple la condición de pendiente
        ($y_1 < y_2$, es decir, la pendiente de la superficie del agua es menor
        que la pendiente de fondo) o la condición de Froude aguas arriba
        ($\mathrm{Fr}_1 \ge 1$) (según la opción \texttt{NORMAL\_FLOW\_LIMITED}),
        entonces
        \begin{equation}
          q \leftarrow \min\!\left(q,\ \beta\,A_1\,R_1^{2/3}\right),
        \end{equation}
        el caudal normal de Manning al tirante aguas arriba.
  \item \textbf{Subrelajación} (iteraciones $> 0$):
        $q = (1-\omega)q_{\mathrm{last}} + \omega\,q$ con acotación del cambio
        de signo a $\pm 0.001$.
  \item \textbf{Límite de caudal del usuario}: $\lvert q \rvert \le
        q_{\mathrm{limit}}$.
  \item \textbf{Compuertas de retención}: el flujo inverso a través de una
        compuerta de retención se anula.
  \item \textbf{Verificación de nodo seco}: el caudal se fuerza a
        $\pm \mathrm{FUDGE}$ si el nodo aguas arriba (aguas abajo) está seco.
\end{enumerate}
Finalmente,
\begin{equation}
  Q^{t+\Delta t} = q\,\mathrm{barrels}, \qquad
  d_{\mathrm{link}} = \min(\bar y, y_{\mathrm{full}}), \qquad
  V_{\mathrm{link}} = \frac{A_1 + A_2}{2}\,L_{\mathrm{cruda}}\,\mathrm{barrels}.
\end{equation}

\subsection{Impulsiones (force mains)}

Las impulsiones (sección \texttt{FORCE\_MAIN}, siempre a sección llena) usan un
núcleo separado (\texttt{process\-Force\-Main\-Link}) con $\sigma = 0$ y
fricción según la ecuación de Hazen--Williams o de Darcy--Weisbach,
\begin{equation}
  \begin{aligned}
    dq_1 &= \Delta t\, g\, \frac{S_f}{\lvert v \rvert},\\
    q &= \frac{q_{\mathrm{old}} - dq_2 + dq_6}{1 + dq_1 + dq_5}.
  \end{aligned}
\end{equation}

% =============================================================================
% 10. CONTINUIDAD EN LOS NODOS, ANEGAMIENTO Y CARGA
% =============================================================================
\section{Continuidad en los nodos, anegamiento y escurrimiento en carga}
\label{sec:node-cont}

\subsection{Caudal neto y cambio de volumen}

En cada nodo, por iteración de Picard (\texttt{DynamicWave.cpp}, función
\texttt{setNodeDepth}), el caudal neto y el cambio de volumen trapezoidal son:
\begin{equation}
  \Delta Q = Q_{\mathrm{in}} - Q_{\mathrm{out}}, \qquad
  \Delta V = \tfrac{1}{2}\left(\Delta Q_{\mathrm{net}}^{t-1} + \Delta Q\right)
  \Delta t,
\end{equation}
donde $\Delta Q_{\mathrm{net}}^{t-1}$ es el caudal neto del paso anterior
(promediado trapezoidal / Crank--Nicolson).

\subsection{Actualización de superficie libre (continuidad EXPLICITA)}

Para un nodo sin escurrimiento en carga, la actualización de tirante es
\begin{equation}
  \Delta y = \frac{\Delta V}{A_S}, \qquad
  y^{t+\Delta t} = y^{t} + \Delta y,
\end{equation}
donde el área superficial del ensamble está pisada por el área superficial
mínima,
\begin{equation}
  A_S = \max\left(A_{SN} + \sum A_{SL},\ A_{\min}\right),
  \qquad A_{\min} = 12.566\ \sfft,
\end{equation}
el valor por defecto (sobrescribible mediante la opción
\texttt{MIN\_SURFAREA}). El área mínima es un dispositivo puramente
computacional; no agrega volumen.

\subsection{Actualización en carga (EXTRAN)}

Un nodo está en carga cuando todos los conductos conectados están a sección
llena o su carga supera la clave de su conducto más alto ($y > y_{\mathrm{crown}}$),
y no es un nodo estanque ni emisario y no está encharcado. Bajo EXTRAN la
ecuación de continuidad se convierte en $\sum Q = 0$, impuesta mediante una
forma de perturbación (Newton/Hardy Cross). La actualización de tirante en
carga es
\begin{equation}
  \Delta y = \frac{\alpha\,\Delta Q}{\mathrm{denom}}, \qquad
  y^{t+\Delta t} = y^{t} + \Delta y,
\end{equation}
con la corrección de nodo terminal $\alpha = 0.6$ para nodos terminales aguas
arriba (solo enlaces de salida, $\mathrm{deg} < 0$) y 1 en caso contrario, y una
mezcla de proximidad a la clave para una transición suave de entrada/salida de
la condición en carga:
\begin{equation}
  \mathrm{denom} = \sum \frac{dQ}{dH} +
  \left(\frac{A_S^{\mathrm{old}}}{\Delta t} - \sum \frac{dQ}{dH}\right)
  e^{-15\,f},
  \qquad f = \frac{y^{t} - y_{\mathrm{crown}}}{y_{\mathrm{crown}}},
\end{equation}
aplicada mientras $y^{t} < 1.25\,y_{\mathrm{crown}}$. La carga se mantiene en o
sobre la clave ($y \ge y_{\mathrm{crown}} - \mathrm{FUDGE}$). El $\sum dQ/dH$
acumulado en cada nodo es la suma de los gradientes por enlace
\eqref{eq:dqdh}; nótese que el motor acumula $\sum dQ/dH$ como una magnitud
positiva con $dQ_{\mathrm{net}}/dH = -\sum dQ/dH$ (una subida de carga aumenta
la salida neta).

\subsection{Continuidad de nodo semi-implícita}

La opción \texttt{NODE\_CONTINUITY} (por defecto EXPLICIT) puede seleccionar la
formulación unificada \emph{semi-implícita}, una única expresión suave válida
tanto para el estado de superficie libre como para el estado en carga.
Linealizando el caudal neto alrededor de la estimación actual de carga con los
gradientes de caudal, $\sum Q^{t+\Delta t} \approx \sum Q + \sum(dQ/dH)\,\Delta H$,
y sustituyendo en la actualización trapezoidal de carga se obtiene
\begin{equation}
  \Delta y = \frac{\Delta V}
    {A_S + \tfrac{\Delta t}{2}\sum \frac{dQ}{dH}}, \qquad
  y^{t+\Delta t} = y^{t} + \Delta y,
\end{equation}
con el denominador pisado por $A_{\min}$. El término de gradiente de caudal
amortigua la actualización (una subida de carga que aumenta la salida neta
reduce el llenado neto), y la expresión es $C^1$-suave a través de la clave,
que es lo que la hace compatible con la aceleración de Anderson.

\subsection{Anegamiento y encharcamiento}

El tirante máximo sin anegar es $y_{\max} = d_{\mathrm{full}} + d_{\mathrm{sur}}$
(para nodos sin encharcamiento; los nodos encharcados nunca se acotan). Si el
tirante candidato supera $y_{\max}$:
\begin{itemize}[nosep]
  \item \textbf{Sin encharcamiento}: la carga se acota a $y_{\max}$ y el exceso
        se pierde como rebalse,
        \begin{equation}
          Q_{\mathrm{ovfl}} = \frac{\Delta V}{\Delta t}, \qquad
          V = V_{\mathrm{full}};
        \end{equation}
  \item \textbf{Con encharcamiento} (con \texttt{ALLOW\_PONDING} y un área de
        encharcamiento positiva): el nodo encharcado actúa como un estanque de
        área constante sobre el brocal,
        \begin{equation}
          V = \max\left(V^{\mathrm{old}} + \Delta V,\ V_{\mathrm{full}}\right), \qquad
          Q_{\mathrm{ovfl}} = \frac{V - \max(V^{\mathrm{old}},\ V_{\mathrm{full}})}
            {\Delta t},
        \end{equation}
        y la carga puede subir libremente sobre el brocal.
\end{itemize}
Un nodo encharcado no puede bajar mucho bajo la profundidad total una vez
encharcado ($y \ge d_{\mathrm{full}} - \mathrm{FUDGE}$).

% =============================================================================
% 11. MÉTODOS DE ESCURRIMIENTO EN CARGA
% =============================================================================
\section{Métodos de escurrimiento en carga}

La opción \texttt{SURCHARGE\_METHOD} selecciona cómo se trata el escurrimiento
presurizado en los conductos cerrados. Los tres métodos actualizan el área
hidráulica de un conducto en carga (tirante $> y_{\mathrm{full}}$) agregando
una ranura angosta de ancho $w_s$:
\begin{equation}
  A = A_{\mathrm{full}} + (y - y_{\mathrm{full}})\,w_s,
\end{equation}
mientras que el radio hidráulico se fija en su valor de sección llena
$R_{\mathrm{full}}$ (la ranura aporta almacenamiento pero no fricción).

\subsection{EXTRAN (por defecto)}

El método de perturbación clásico de la Sección~\ref{sec:node-cont}:
actualizaciones de nodo en carga basadas en $dQ/dH$ con un corte de clave en
$y/y_{\mathrm{full}} = 0.96$ usado para el ancho empleado en el cálculo del
área superficial.

\subsection{Ranura de Preissmann estática}

El conducto lleva una ranura ficticia angosta sobre la clave para que la
formulación de superficie libre siga siendo válida durante toda la
presurización. El ancho de ranura sigue la fórmula de Sj{\"o}berg (1982),
\begin{equation}
  \frac{w_s}{W_{\max}} = 0.5423\,
  \exp\!\left(-\left(\frac{y}{y_{\mathrm{full}}}\right)^{2.4}\right),
\end{equation}
aplicada para $0.985257 \le y/y_{\mathrm{full}} \le 1.78$, fijada en
$w_s = 0.01\,W_{\max}$ sobre ese rango y cero bajo el corte.

\subsection{Ranura de Preissmann dinámica}

Una extensión de OpenSWMM (Sharior, Hodges \& Vasconcelos 2023) en la que el
área de la ranura evoluciona en el tiempo como un elemento de almacenamiento
transitorio. El ancho superior de la ranura está gobernado por una celeridad de
onda de presión objetivo $c_{pT}$ y por un número de Preissmann variable en el
tiempo $P$:
\begin{equation}
  T_s = \frac{g\,A_{\mathrm{full}}}{c_{pT}^{2}}\,P^{2},
  \qquad
  P_{0} = \max\!\left(\frac{c_{pT}}{\alpha\,c_g},\ 1\right),
  \qquad c_g = \sqrt{g\,A_{\mathrm{full}}/W_{\max}},
\end{equation}
con el parámetro de choque $\alpha = 3$ (por defecto). El área de ranura
acumulada depende de la trayectoria,
\begin{equation}
  A_s \leftarrow \max\!\left(A_s + T_s\,\Delta h_s,\ 0\right),
  \qquad h_s = \max(\bar y - y_{\mathrm{full}},\ 0),
\end{equation}
y $P$ decae exponencialmente después del inicio de la presurización y se suaviza
espacialmente entre los nodos. Mientras una ranura está activa, el ancho
superior es $T_s$ y el área superficial aportada a un nodo extremo en carga es
$T_s\,L/4$. Las cargas de los nodos se actualizan en todo momento con la
fórmula ordinaria de superficie libre; la rama de carga nunca se invoca. El
paso de tiempo variable usa la celeridad de presión $c_p = c_{pT}/P$.

% =============================================================================
% 12. ESTRUCTURAS HIDRÁULICAS NO-CONDUCTO
% =============================================================================
\section{Estructuras hidráulicas no-conducto}
\label{sec:structures}

Los caudales de las estructuras no-conducto los calcula \texttt{HydStructures.cpp}
(\texttt{StructureSolver}) dentro de la iteración de Picard de la onda
dinámica, un enlace a la vez en orden de índice (para que dos bombas que
comparten una cámara húmeda se vean mutuamente la extracción), con distribución
inmediata en los acumuladores de los nodos y subrelajación a $\omega = 0.5$
para iteraciones $> 0$ (las bombas quedan exentas).

\subsection{Bombas}

El control de arranque/parada de las bombas usa histéresis de tirante aplicada
una vez por paso de tránsito:
\begin{equation}
  \text{detener si } d < d_{\mathrm{off}} \ (\text{regulación} > 0),
  \qquad
  \text{arrancar si } d > d_{\mathrm{on}} \ (\text{regulación} = 0).
\end{equation}
Los tipos de curva de bomba (1--5) determinan el caudal $q$ a partir de la curva
versus el volumen de la cámara húmeda (tipo 1), el tirante (tipos 2 y 4) o la
carga (tipos 3 y 5, donde $q$ se escala por la regulación de velocidad $s$ y la
carga por $s^2$); una bomba ideal descarga exactamente el ingreso más el rebalse
del nodo aguas arriba. La limitación de caudal (legado
\texttt{getModPumpFlow}) evita bombear más de lo disponible (protección contra
aspiración en seco) y, para bombas tipo 1 o alimentadas por estanque, acota $q$
por $Q_{\mathrm{in}} + V_{\mathrm{old}}/\Delta t$.

\subsection{Orificios}

Los orificios de fondo y laterales usan un coeficiente y un tirante crítico
definidos a partir de la altura de apertura $h_{\mathrm{open}} =
s\,y_{\mathrm{full}}$ (regulación $s$):
\begin{align}
  f &= \min\!\left(\frac{\text{carga}}{h_{\mathrm{crit}}},\ 1\right),
  \qquad h_{\mathrm{crit}} = \frac{C_d}{0.414}\,\frac{h_{\mathrm{open}}}
    {4}\ \text{(circular)}\ \text{o}\ \frac{C_d}{0.414}\,
    \frac{h_{\mathrm{open}}\,W_{\max}}{2(h_{\mathrm{open}} + W_{\max})}
    \text{(rectangular)},\\
  q &= \begin{cases}
    C_w\,f^{3/2} & f < 1 \quad\text{(escurrimiento parcial tipo vertedero)},\\
    C_d\,A_{\mathrm{eff}}\sqrt{2g\,H} & f \ge 1 \quad\text{(orificio lleno)},
  \end{cases}
\end{align}
con $C_w = C_d\sqrt{h_{\mathrm{crit}}}\,A_{\mathrm{eff}}\sqrt{2g}$
(coeficiente de vertedero de cresta afilada 0.414), $A_{\mathrm{eff}}$ el área
de la sección transversal a la altura de apertura y $H$ la carga apropiada
(carga diferencial cuando está sumergido/ahogado). El gradiente de caudal es
\begin{equation}
  \frac{dQ}{dH} = \frac{3}{2}\,\frac{q}{f\,h_{\mathrm{crit}}}
  \quad\text{(régimen vertedero)}, \qquad
  \frac{dQ}{dH} = \frac{q}{2H}
  \quad\text{(régimen orificio)}.
\end{equation}
Puede aplicarse una pérdida de carga opcional de compuerta de retención ARMCO y
la corrección de sumersión de Villemonte.

\subsection{Vertederos}

Los cuatro tipos de vertedero (transversal, lateral, escotadura en V,
trapezoidal) descargan según
\begin{equation}
  \begin{aligned}
    q &= C_d\,L\,h^{3/2} &&\text{(transversal)},\\
    q &= C_d\,L^{0.83}\,h^{1.67} &&\text{(lateral, corregido)},\\
    q &= C_d\,s_h\,h^{5/2} &&\text{(escotadura en V)},
  \end{aligned}
\end{equation}
con contracciones de extremo que reducen la longitud efectiva. Cuando la línea
de gradiente hidráulica aguas arriba alcanza la clave ($H_1 \ge
H_{\mathrm{crown}}$) y se permite la entrada en carga, el vertedero transita a
escurrimiento por orificio usando un coeficiente equivalente de orificio
calculado desde el caudal del vertedero a apertura completa:
\begin{equation}
  C_{\mathrm{sur}} = \frac{Q_{\mathrm{weir}}(s\,y_{\mathrm{full}})}
    {\sqrt{(s\,y_{\mathrm{full}})/2}},
  \qquad
  q = C_{\mathrm{sur}}\sqrt{H_{\mathrm{orif}}},
\end{equation}
con $H_{\mathrm{orif}}$ la carga al punto medio de la abertura (o la carga
diferencial cuando está sumergido/ahogado). Si no se permite la entrada en
carga, la carga se acota en la clave y se mantiene la ecuación de vertedero. La
corrección de sumersión de Villemonte se aplica cuando la LGH aguas abajo
supera la cresta. La contribución de área superficial del vertedero se anula
por compatibilidad con SWMM 4.

\subsection{Descargas (outlets)}

Las descargas usan curvas de gasto en función de la carga o del tirante
(tabulares o funcionales):
\begin{equation}
  q = C\,H^{e} \quad\text{(funcional)}, \qquad
  q = \mathrm{table}(H) \quad\text{(tabular)},
\end{equation}
escaladas por la regulación. El $dQ/dH$ se deja deliberadamente en cero
(igual que el legado), lo que evita inflar el denominador de carga de los
nodos.

% =============================================================================
% 13. CONDICIONES DE BORDE DE EMISARIO
% =============================================================================
\section{Condiciones de borde de emisario}
\label{sec:outfall}

Los niveles de agua de los emisarios se fijan al inicio de cada paso de
tránsito y nuevamente dentro de cada iteración de Picard a partir de los
caudales actuales de los conductos. Para el conducto conectado a un emisario en
un extremo de elevación $z$, con caudal por barril $q$, se calculan el tirante
normal $y_n$ (desde el factor de sección inverso) y el tirante crítico $y_c$
($y_c$ tiene formas cerradas para las formas estándar y una búsqueda numérica
de raíz en el resto). El nivel depende del tipo de condición de borde
(\texttt{Outfall.cpp:setOutfallDepth}):
\begin{itemize}[nosep]
  \item \textbf{LIBRE (FREE)}: caída libre, $d = z + \min(y_n, y_c)$.
  \item \textbf{NORMAL}: $d = z + y_n$.
  \item \textbf{FIJO (FIXED)}, \textbf{MAREA (TIDAL)}, \textbf{SERIE DE TIEMPO
        (TIMESERIES)}: el nivel $s_{\mathrm{stage}}$ es constante, proviene de
        una tabla de mareas (hora del día) o de una serie de tiempo; entonces
        \begin{equation}
          d = \begin{cases}
            s_{\mathrm{stage}} - z_{\mathrm{inv}} & z + y_c + z_{\mathrm{inv}}
              < s_{\mathrm{stage}},\\
            \max(0,\ s_{\mathrm{stage}} - z_{\mathrm{inv}}) & z > 0 \wedge
              s_{\mathrm{stage}} < z_{\mathrm{inv}} + z,\\
            z + y_c & z > 0 \wedge s_{\mathrm{stage}} \ge z_{\mathrm{inv}} + z,\\
            y_c & z = 0.
          \end{cases}
        \end{equation}
\end{itemize}
El último caso mantiene un emisario fijo de descarga libre al tirante crítico
del conducto. Una compuerta de retención en un emisario bloquea el flujo
inverso.

% =============================================================================
% 14. DIVISORES DE FLUJO
% =============================================================================
\section{Divisores de flujo}

Los divisores están activos solo bajo tránsito por onda cinemática (y
permanente); bajo onda dinámica actúan como cámaras de unión ordinarias. Dado
el ingreso total $Q_{\mathrm{in}}$ en el nodo divisor:
\begin{itemize}[nosep]
  \item \textbf{De corte (cutoff)}: desvía todo el ingreso sobre un mínimo,
        $Q_{\mathrm{div}} = \max(0, Q_{\mathrm{in}} - q_{\min})$.
  \item \textbf{De rebalse (overflow)}: desvía el ingreso sobre la capacidad del
        enlace de continuación.
  \item \textbf{Tabular}: $Q_{\mathrm{div}} = Q_{\mathrm{in}}\cdot
        f(Q_{\mathrm{in}})$ con la fracción de una tabla de gasto.
  \item \textbf{Vertedero}: $Q_{\mathrm{div}} = C_d\,W\,d^{3/2}$ a partir de
        una relación de vertedero sobre el tirante del nodo, acotada al ingreso.
\end{itemize}
El caudal desviado se escribe en el enlace de desvío; el enlace de continuación
lleva el resto.

% =============================================================================
% 15. PASO DE TIEMPO ADAPTATIVO Y CONDICIÓN DE COURANT
% =============================================================================
\section{Paso de tiempo adaptativo y condición de Courant}

\subsection{Condición de Courant--Friedrichs--Lewy}

Como el esquema de onda dinámica actualiza caudales y cargas elemento por
elemento (no hay acoplamiento espacial simultáneo), está sujeto a la condición
de Courant: el paso de tiempo no debe exceder el tiempo que tarda una onda
dinámica en atravesar el conducto más corto,
\begin{equation}
  \Delta t \le \frac{L}{\lvert U \rvert + c},
  \qquad c = \sqrt{g\,\frac{A}{W}}
  \text{ (celeridad de la onda de gravedad).}
\end{equation}

\subsection{Paso basado en el enlace}

Para cada conducto (\texttt{getLinkStep}, \texttt{Dynamic\-Wave.cpp:3711})
con caudal por barril $q = \lvert Q \rvert/\mathrm{barrels}$, área media $A$ y
número de Froude $\mathrm{Fr}$:
\begin{equation}
  t = \frac{V/\mathrm{barrels}}{q}\,
      \frac{L'}{L}\,\frac{\mathrm{Fr}}{1 + \mathrm{Fr}},
\end{equation}
(la expresión CFL clásica de SWMM, escalada por la razón longitud
alargada/cruda). Bajo la ranura de Preissmann dinámica, un conducto en carga
usa en su lugar la celeridad de presión,
\begin{equation}
  t = \frac{L\,(L'/L)}{\lvert v \rvert + c_p}, \qquad
  c_p = \frac{c_{pT}}{P}.
\end{equation}

\subsection{Paso basado en el nodo}

Para cada nodo no-emisario bajo la clave con tasa de cambio de tirante no nula
$\dot y = \lvert \Delta y/\Delta t \rvert$:
\begin{equation}
  t = 0.25\,\frac{y_{\mathrm{crown}}}{\dot y},
\end{equation}
resguardando contra un cambio excesivo de carga en un solo paso.

\subsection{Combinación de las restricciones}

El siguiente paso de tránsito (\texttt{getRoutingStep}) es el mínimo sobre
todos los enlaces, nodos y pares de unión virtual, cada uno escalado por el
factor de Courant del usuario, y luego pisado y acotado:
\begin{equation}
  \Delta t_{\min}^{\mathrm{eff}} =
    \max\!\left(\min_{j,i}\{t_j, t_i\}\cdot \mathrm{Cr},\
      \max(\min(\Delta t_{\min},\ \Delta t_{\text{tránsito}}),\
      0.001\ \st)\right),
\end{equation}
redondeado a milisegundos. El paso efectivo se acota luego por el
\texttt{ROUTING\_STEP} fijo del usuario y por la duración restante de la
simulación. El paso inicial es el paso mínimo. La primera llamada (sin caudales
aún) también retorna el paso mínimo. El solver FV sub-pasea internamente a su
propio límite CFL, por lo que bajo FV el paso de tránsito es solo una cadencia
de reporte.

% =============================================================================
% 16. ESTABILIDAD, CONVERGENCIA Y BALANCE DE MASA
% =============================================================================
\section{Estabilidad, convergencia y balance de masa}

\subsection{Indicadores de estabilidad}

La inestabilidad numérica se manifiesta como oscilaciones que no se amortiguan
en el caudal y en la superficie del agua, y en nodos que se secan repetidamente.
Dos métricas de reporte la cuantifican:
\begin{itemize}[nosep]
  \item el \textbf{error de continuidad de caudal} global (abajo), y
  \item el \textbf{índice de inestabilidad de flujo del enlace (FII)} --- el
        conteo normalizado de las veces que el caudal de un enlace supera a sus
        dos vecinos.
\end{itemize}

\subsection{Convergencia del ciclo de Picard}

Un paso de tránsito ``no converge'' solo cuando la tolerancia de carga
($\varepsilon_H = 0.005\ \ft$) no se cumple en algún nodo no-emisario después de
\texttt{MAX\_TRIALS} iteraciones. La aceleración de Anderson (opcional) acelera
la convergencia mezclando las dos salidas más recientes del operador en cada
nodo:
\begin{equation}
  \alpha_k = \mathrm{clamp}\!\left(
    \frac{r_k\,(r_k - r_{k-1})}{(r_k - r_{k-1})^{2}},\ 0,\ 1\right),
  \qquad
  H_{k+1} = (1-\alpha_k)\,G(H_k) + \alpha_k\,G(H_{k-1}),
\end{equation}
con $r_k = G(H_k) - H_k$ el residual y $G$ el operador de actualización de
carga. Se omite en nodos donde $G$ es no suave (en carga bajo EXTRAN, ranura
dinámica activa, cerca del corte de la ranura estática, vertedero/orificio en
la clave, extremos de bomba, borde de encharcamiento).

\subsection{Balance de masa del tránsito}

El error de continuidad del tránsito (fracción del ingreso total) es
\begin{equation}
  \varepsilon_{\mathrm{routing}} = \frac{Q_{\mathrm{dw}} + Q_{\mathrm{wet}} +
    Q_{\mathrm{gw}} + Q_{\mathrm{rdii}} + Q_{\mathrm{ext}} +
    V_{\mathrm{init}} - \left(Q_{\mathrm{aneg}} + Q_{\mathrm{sal}} +
    Q_{\mathrm{evap}} + Q_{\mathrm{filt}} + V_{\mathrm{final}}\right)}
    {Q_{\mathrm{total\,in}}},
\end{equation}
donde los ingresos son de tiempo seco, tiempo húmedo, napa freática, RDII y
externos, y las salidas son anegamiento, caudal de salida del sistema,
evaporación, filtración y almacenamiento final. El anegamiento solo se cuenta
cuando el volumen del nodo no excede su volumen total; la descarga de los
emisarios cuenta como caudal de salida del sistema, y el flujo inverso en los
emisarios cuenta como ingreso externo. Las pérdidas por evaporación/filtración
de conductos y estanques también se cargan aquí.

% =============================================================================
% 17. OPCIONES Y VALORES POR DEFECTO
% =============================================================================
\section{Opciones y valores por defecto}

La Tabla~\ref{tab:options} resume las opciones relacionadas con el tránsito
(\texttt{SimulationOptions.hpp}) y sus valores por defecto.

\begin{longtable}{>{\raggedright}p{4.6cm}>{\raggedright}p{2.6cm}>{\raggedright\arraybackslash}p{7.0cm}}
\toprule
\textbf{Opción} & \textbf{Defecto} & \textbf{Propósito} \\
\midrule
\endhead
\texttt{FLOW\_ROUTING} & DYNWAVE & formulación de tránsito \\
\texttt{ROUTING\_STEP} & 20 s & paso de tránsito fijo \\
\texttt{MINIMUM\_STEP} & 0.5 s & piso del paso CFL \\
\texttt{DRY\_STEP} / \texttt{WET\_STEP} & 3600 / 300 s & límites del paso de escorrentía \\
\texttt{VARIABLE\_STEP} & 0.75 & factor de Courant (0 = paso fijo) \\
\texttt{LENGTHENING\_STEP} & 0 & paso de alargamiento de Courant \\
\texttt{MAX\_TRIALS} & 8 & límite de intentos de Picard \\
\texttt{HEAD\_TOLERANCE} & 0.005 (unidades del proyecto) & tolerancia de convergencia de nodo \\
\texttt{SURCHARGE\_METHOD} & EXTRAN & EXTRAN / SLOT / DYNAMIC\_SLOT \\
\texttt{NODE\_CONTINUITY} & EXPLICIT & dos ramas vs.\ semi-implícita \\
\texttt{INERTIAL\_DAMPING} & PARTIAL & política de amortiguación por Froude \\
\texttt{NORMAL\_FLOW\_LIMITED} & BOTH & tope de caudal normal por pendiente/Froude \\
\texttt{MIN\_SURFAREA} & 12.566 ft$^2$ & piso del área de nodo \\
\texttt{ALLOW\_PONDING} & false & permite encharcamiento sobre el brocal \\
\texttt{ANDERSON\_ACCEL} & false & aceleración de Anderson \\
\texttt{SKIP\_STEADY\_STATE} & false & omite el tránsito en régimen permanente \\
\bottomrule
\caption{Opciones de tránsito y valores por defecto.}
\label{tab:options}
\end{longtable}

% =============================================================================
% 18. REFERENCIAS
% =============================================================================
\section{Referencias}

\begin{enumerate}[nosep,label={[\arabic*]}]
  \item Rossman, L.~A. (2017). \emph{Storm Water Management Model Reference
        Manual Volume II -- Hydraulics}. U.S.\ EPA.
  \item Código fuente del motor OpenSWMM:
        \begin{itemize}[nosep,itemsep=0pt]
          \item \src{hydraulics/Routing.cpp} \\
                \src{hydraulics/DynamicWave.cpp} \\
                \src{hydraulics/KinematicWave.cpp} \\
                \src{hydraulics/HydStructures.cpp}
          \item \src{hydraulics/Node.cpp} \\
                \src{hydraulics/Link.cpp} \\
                \src{hydraulics/Outfall.cpp} \\
                \src{hydraulics/Divider.cpp}
          \item \src{hydraulics/XSectBatch.cpp} \\
                \src{hydraulics/XSectKernels.hpp} \\
                \src{core/SimulationOptions.hpp} \\
                \src{core/SWMMEngine.cpp}
        \end{itemize}
  \item Sharior, S., Hodges, B.~R., \& Vasconcelos, J.~G. (2023).
        Generalized, dynamic, and transient-storage form of the Preissmann
        slot. \emph{Journal of Hydraulic Engineering}, 149(11).
  \item Sj{\"o}berg, A. (1982). Sewer network models dagnum and DIVISION --
        a brief description. \emph{Swedish Water and Wastewater Works
        Association}.
  \item Roesner, L.~A., Nichandros, H.~M., Shubinski, R.~P., Feldman, A.~D.,
        Abbott, J.~W., \& Delleur, J.~W. (1981). \emph{Physical model for
        combined sewer overflows}. Proc.\ ASCE.
  \item de~Almeida, G.~A.~M., \& Bates, P.~D. (2013). Applicability of the
        local inertial approximation of the shallow water equations to flood
        modeling. \emph{Water Resources Research}, 49(8).
\end{enumerate}

\end{document}
